\input{../tex-inputs/latex-leseansicht-vorspann}

\section[Rainer Maria Rilke an Arthur Schnitzler, {[}zwischen 5. und 16.? 5. 1901{]}]{L04538 Rainer Maria Rilke an Arthur Schnitzler, {[}zwischen 5. und 16.? 5. 1901{]}}
\nopagebreak\mylabel{L04538v}
\rehead{ }\normalsize\beginnumbering\briefempfaengerindex{Schnitzler, Arthur@\textsc{Schnitzler, Arthur}!zzzRilke, Rainer Maria@\emph{von Rainer Maria Rilke}!1901-05-161@{{[}zwischen 5. und 16.? 5. 1901{]}}|(be}
\toendnotes[C]{\smallbreak\pagebreak[2]}
\correspDesc{Versand  durch Rainer Maria Rilke im Zeitraum [zwischen 5. und
                  16.? 5. 1901] in Dresden
\newline{}Erhalt  durch Arthur Schnitzler in Wien}\toendnotes[C]{\smallbreak}
\Standort{CUL, Schnitzler, B 84.}
\physDesc{Brief, 1 Blatt, 3 Seiten, 2879 Zeichen
\newline{}Handschrift: schwarze Tinte, deutsche Kurrent
\newline{}Schnitzler: mit rotem Buntstift eine Unterstreichung }
\buchAbdrucke{\weitereDrucke{\emph{Rainer Maria Rilke und Arthur Schnitzler. Ihr Briefwechsel.}. Mit Anmerkungen herausgegeben von Heinrich Schnitzler In: \emph{Wort und Wahrheit}, Jg. 13, Nr. 4, April 1958, S. 282–298, hier S. 285–286.} }\toendnotes[C]{\smallbreak}
\pstart
           \raggedleft{}{\pb}\textsc{Dresden, Weißer Hirsch\oindex{Weißer Hirsch@\textbf{Weißer Hirsch}|pw},}\pend
           
\pstart
           \raggedleft{}\textsc{Sanatorium D\textsuperscript{r}
                        Lahmann\oindex{Physiatrisches Sanatorium (Lahmann)@\textbf{Physiatrisches Sanatorium (Lahmann)}|pw}.}\pend
           
\pstart{}Lieber und verehrter Doctor \textsc{Schnitzler},\pend\vspace{0.5em}
\pstart
           Ihr Buch »\textsc{Der Schleier der Beatrice\pwindex{Schnitzler, Arthur 15. 5. 1862 Wien – 21. 10. 1931 ebd.@\textsc{Schnitzler, Arthur} (15. 5. 1862 Wien – 21. 10. 1931 ebd.)!Schleier der Beatrice. Schauspiel in fünf Akten@\strich\emph{Der Schleier der Beatrice. Schauspiel in fünf Akten}|pw}}« iſt mir {[}hier{]}her nachgeſendet worden und Ihre lieben
               Wünſche auch. Sie haben mit Ihrem Geſchenk ganz unabſichtlich eine{ }ſo gute Zeit
               getroffen; ich nehme das wie ein frohes Zeichen auf. Ich empfing es \label{K_L04538-1v}\edtext{nichtmehr allein}{\lemma{\textnormal{\emph{nichtmehr allein}}}\Cendnote{\textnormal{Am 28. 4. 1901 hatten Clara Westhoff\pwindex{Westhoff, Clara 21.\,11.\,1878 Bremen – 9.\,3.\,1954 Fischerhude@\textsc{Westhoff, Clara} (21.\,11.\,1878 Bremen – 9.\,3.\,1954 Fischerhude)|pwk} und Rilke\pwindex{Rilke, Rainer Maria 4.\,12.\,1875 Prag – 29.\,12.\,1926 Montreux@\textsc{Rilke, Rainer Maria} (4.\,12.\,1875 Prag – 29.\,12.\,1926 Montreux)|pwk}{ }geheiratet\eventindex{Bremen@\textbf{Bremen}!Hochzeit von Clara Westhoff und Rainer Maria Rilke, 28.4.1901@Hochzeit von Clara Westhoff und Rainer Maria Rilke, 28.4.1901|pwkv}. }}}\label{K_L04538-1} und
               doppelt danke ich Ihnen deshalb dafür.\pend
           
\pstart
           Ich habe das Drama\pwindex{Schnitzler, Arthur 15. 5. 1862 Wien – 21. 10. 1931 ebd.@\textsc{Schnitzler, Arthur} (15. 5. 1862 Wien – 21. 10. 1931 ebd.)!Schleier der Beatrice. Schauspiel in fünf Akten@\strich\emph{Der Schleier der Beatrice. Schauspiel in fünf Akten}|pwv} zuerst
               allein geleſen, ehe ich es meiner Frau\pwindex{Westhoff, Clara 21.\,11.\,1878 Bremen – 9.\,3.\,1954 Fischerhude@\textsc{Westhoff, Clara} (21.\,11.\,1878 Bremen – 9.\,3.\,1954 Fischerhude)|pwv} vorlas. Dann haben wir viel darüber geſprochen, und{ }ſind in beſtimmten letzten Punkten ganz einig geweſen.\pend
           
\pstart
           Ich danke Ihnen gerade für dieſes Geſchenk; denn mir liegt »der \textsc{Schleier der Beatrice\pwindex{Schnitzler, Arthur 15. 5. 1862 Wien – 21. 10. 1931 ebd.@\textsc{Schnitzler, Arthur} (15. 5. 1862 Wien – 21. 10. 1931 ebd.)!Schleier der Beatrice. Schauspiel in fünf Akten@\strich\emph{Der Schleier der Beatrice. Schauspiel in fünf Akten}|pw}}« in der aufſteigenden Linie derjenigen Ihrer Werke, die ich am meiſten liebe.
                  »\textsc{Liebelei\pwindex{Schnitzler, Arthur 15. 5. 1862 Wien – 21. 10. 1931 ebd.@\textsc{Schnitzler, Arthur} (15. 5. 1862 Wien – 21. 10. 1931 ebd.)!Liebelei. Schauspiel in drei Akten@\strich\emph{Liebelei. Schauspiel in drei Akten}|pw}}« und »\textsc{Die Gefährtin\pwindex{Schnitzler, Arthur 15. 5. 1862 Wien – 21. 10. 1931 ebd.@\textsc{Schnitzler, Arthur} (15. 5. 1862 Wien – 21. 10. 1931 ebd.)!Gefährtin. Schauspiel in einem Akt@\strich\emph{Die Gefährtin. Schauspiel in einem Akt}|pw}}«{ }ſind direkte Vorfahren dieſes Buches – für mein Gefühl. – Wie in der Gefährtin\pwindex{Schnitzler, Arthur 15. 5. 1862 Wien – 21. 10. 1931 ebd.@\textsc{Schnitzler, Arthur} (15. 5. 1862 Wien – 21. 10. 1931 ebd.)!Gefährtin. Schauspiel in einem Akt@\strich\emph{Die Gefährtin. Schauspiel in einem Akt}|pw} die Geſtalt einer Toten {\pb}mit einer gewiſſen Rückſichtsloſigkeit{ }ſich ausbaut und vollendet, bis der letzte Zweifel, der sie ungewiſs macht in den
               Umriſſen, fällt —{ }ſo flackert hier das Leben einer Sterbenden\pwindex{Schnitzler, Arthur 15. 5. 1862 Wien – 21. 10. 1931 ebd.@\textsc{Schnitzler, Arthur} (15. 5. 1862 Wien – 21. 10. 1931 ebd.)!Schleier der Beatrice. Schauspiel in fünf Akten@\strich\emph{Der Schleier der Beatrice. Schauspiel in fünf Akten}|pwv} über den vielen Geſtalten einer bewegten und
               wilden Zeit. Wunderbar einfach iſt das komplizierte Weſen der \textsc{Beatrice\pwindex{Schnitzler, Arthur 15. 5. 1862 Wien – 21. 10. 1931 ebd.@\textsc{Schnitzler, Arthur} (15. 5. 1862 Wien – 21. 10. 1931 ebd.)!Schleier der Beatrice. Schauspiel in fünf Akten@\strich\emph{Der Schleier der Beatrice. Schauspiel in fünf Akten}|pwv}} gefaſst und mit wirklich großer Gerechtigkeit{ }ſtehen Sie über ihm und{ }ſeinen
               Wirren. So daſs es{ }ſcheint, als{ }ſtünde man gar keinem Ausnahmsfall, im Gegentheil
               einem täglichen (bislang unerkannten)\textsc{Typus} gegenüber; es
               iſt faſt, als wäre nur für Miſstrauiſche und Schwerfällige das fantaſtische\textsc{Milieu} jener breiten Zeit nothwendig geweſen, und man
               bedauert einen Augenblick faſt, die Geſtalt des bewegten Mädchen\pwindex{Schnitzler, Arthur 15. 5. 1862 Wien – 21. 10. 1931 ebd.@\textsc{Schnitzler, Arthur} (15. 5. 1862 Wien – 21. 10. 1931 ebd.)!Schleier der Beatrice. Schauspiel in fünf Akten@\strich\emph{Der Schleier der Beatrice. Schauspiel in fünf Akten}|pwv} aus{ }ſo vielen Gruppen auslöſen zu
               müſſen. Aber das Bedauern geht vorüber: \textsc{Beatrice\pwindex{Schnitzler, Arthur 15. 5. 1862 Wien – 21. 10. 1931 ebd.@\textsc{Schnitzler, Arthur} (15. 5. 1862 Wien – 21. 10. 1931 ebd.)!Schleier der Beatrice. Schauspiel in fünf Akten@\strich\emph{Der Schleier der Beatrice. Schauspiel in fünf Akten}|pwv}} prägt{ }ſich{ }ſo{ }ſtark aus, daſs die ganze Zeit mit \strikeout{Ihre} ihren Kriegen und Giften, Liedern und Schreien, hingeriſſen von \uline{ihr},{ }ſich{ }ſo zu geberden{ }ſcheint.\pend
           
\pstart
           Ich möchte Ihnen noch viel{ }ſagen zu diesem Buche\pwindex{Schnitzler, Arthur 15. 5. 1862 Wien – 21. 10. 1931 ebd.@\textsc{Schnitzler, Arthur} (15. 5. 1862 Wien – 21. 10. 1931 ebd.)!Schleier der Beatrice. Schauspiel in fünf Akten@\strich\emph{Der Schleier der Beatrice. Schauspiel in fünf Akten}|pwv}: z. B. auch daſs »das Herzogin{\pb}seinwollen« \textsc{Beatrice’s\pwindex{Schnitzler, Arthur 15. 5. 1862 Wien – 21. 10. 1931 ebd.@\textsc{Schnitzler, Arthur} (15. 5. 1862 Wien – 21. 10. 1931 ebd.)!Schleier der Beatrice. Schauspiel in fünf Akten@\strich\emph{Der Schleier der Beatrice. Schauspiel in fünf Akten}|pwv}} (ſchon in der erſten Erzählung des Traumes kommt das zum Ausdruck) – eigentlich
               wie ein Künſtlerwunſch wirkt, die \uline{Stimmung Herzogin}
               irgendwie zu{ }ſchaffen. Ich weiß nicht, ob das klar genug geſagt iſt. Es wäre das
               Schicksal \uline{Nicht}ſchöpferischer, die \uline{doch} Schaffenstriebe tragen;{ }ſie müſſen{ }ſie nothwendig
               miſsverſtehen und die ihnen vorſchwebenden Geſtalten im Leben durchſetzen wollen{\dotsfour}\pend
           
\pstart
           Vielleicht{ }ſag ich das{ }ſpäter einmal einfacher, deutlicher. Ich bin hier jetzt einer{ }ſtrengen \textsc{Cur} unterworfen, der Worte und des Ausdrucks
               entwöhnt {\dots} und dies{ }ſoll nur ein raſcher Dank{ }ſein.\pend
           
\pstart
           Meine Frau\pwindex{Westhoff, Clara 21.\,11.\,1878 Bremen – 9.\,3.\,1954 Fischerhude@\textsc{Westhoff, Clara} (21.\,11.\,1878 Bremen – 9.\,3.\,1954 Fischerhude)|pwv} grüßt Sie{ }ſehr.\pend
           
\pstart
           Übrigens:{ }ſie iſt Bildhauerin. Hat heuer in \textsc{München\oindex{München@\textbf{München}|pw}} und \textsc{Dresden\oindex{Dresden@\textbf{Dresden}|pw}} ausgeſtellt. Vielleicht kommen Sie in die eine oder andere Stadt. Deshalb
               erwähne ichs, da Sie{ }ſich vielleicht dafür intereſſieren.\pend
           
\pstart
           Und vielleicht kommt es auch einmal zu einem Wiederſehen. Unſer Bauernhaus\oindex{?? [Bauernhaus von Clara und Rainer Maria Rilke in Westerwede]@\textbf{?? [Bauernhaus von Clara und Rainer Maria Rilke in Westerwede]}|pwv} iſt in nächſter Nachbarſchaft
               von \textsc{Worpswede\oindex{Worpswede@\textbf{Worpswede}|pw}}. Es verlohnt{ }ſich alſo, einmal zu uns zu kommen.\pend
           
\pstart
           Mit viel lieben Grüßen in größter Ergebenheit{\\[\baselineskip]}Ihr:{\\[\baselineskip]}\spacefill\mbox{RainerMariaRilke}\pend
           \leftskip=0em{}\selectlanguage{ngerman}\endnumbering\briefempfaengerindex{Schnitzler, Arthur@\textsc{Schnitzler, Arthur}!zzzRilke, Rainer Maria@\emph{von Rainer Maria Rilke}!1901-05-051@{{[}zwischen 5. und 16.? 5. 1901{]}}|)be}\mylabel{L04538h}
\begin{anhang}
\end{anhang}\newcommand{\dateiname}{L04538}\newcommand{\titel}{Rainer Maria Rilke an Arthur Schnitzler, [zwischen 5. und 16.? 5. 1901]}\newcommand{\editorInnen}{Herausgegeben von Jahnke, SelmaMüller, Martin Anton}\input{../tex-inputs/latex-leseansicht-abspann}
